\documentclass[a4paper,12pt]{article}

\usepackage[T2A]{fontenc}
\usepackage[utf8]{inputenc}
\usepackage[russian]{babel}
\usepackage{amsmath,amssymb,amsthm,mathtools}
\usepackage{helvet}
\usepackage{PTSans}
\renewcommand{\familydefault}{\sfdefault}
\usepackage{icomma}
\usepackage[a4paper,left=20mm,right=20mm,top=20mm,bottom=22mm]{geometry}
\usepackage{microtype}
\usepackage{tikz}
\usetikzlibrary{calc,arrows.meta,positioning,shapes.geometric}
\usepackage{pgfplots}
\pgfplotsset{compat=1.18}
\usepackage{tabularx,booktabs}
\usepackage{enumitem}
\usepackage{parskip}
\usepackage{fancyhdr}
\usepackage{setspace}
\usepackage{xcolor}

\definecolor{accent}{HTML}{008080}
\definecolor{darkaccent}{HTML}{004D4D}
\definecolor{textgray}{HTML}{333333}
\definecolor{softgray}{HTML}{F3F6F6}

\setlength{\headheight}{25pt}
\setlength{\parindent}{0pt}
\setlength{\parskip}{0.55em}
\onehalfspacing
\setlist[itemize]{leftmargin=1.5em,itemsep=0.25em,topsep=0.3em}
\setlist[enumerate]{leftmargin=1.7em,itemsep=0.4em,topsep=0.3em}
\renewcommand{\arraystretch}{1.35}

\pagestyle{fancy}
\fancyhf{}
\fancyhead[L]{\small\textcolor{textgray}{ЕГЭ. Профильная математика}}
\fancyhead[R]{\small\textcolor{textgray}{Графики показательных функций}}
\fancyfoot[C]{\small\thepage}
\renewcommand{\headrulewidth}{0.3pt}
\renewcommand{\headrule}{%
  \hbox to\headwidth{\color{accent}\leaders\hrule height \headrulewidth\hfill}%
}

\newcommand{\sectionline}{%
  \vspace{-0.2em}
  \textcolor{accent}{\rule{\linewidth}{0.5pt}}
  \vspace{0.25em}
}
\newcommand{\key}[1]{\textbf{\textcolor{darkaccent}{#1}}}

\tikzset{
  flowstart/.style={
    ellipse,
    draw=accent,
    line width=0.9pt,
    align=center,
    minimum width=5cm,
    minimum height=10mm,
    font=\small
  },
  flowstep/.style={
    rounded corners=3pt,
    draw=accent,
    line width=0.9pt,
    align=center,
    text width=10.8cm,
    minimum height=10mm,
    inner sep=5pt,
    font=\small
  },
  flowdecision/.style={
    diamond,
    aspect=3.1,
    draw=accent,
    line width=0.9pt,
    align=center,
    text width=4.1cm,
    inner sep=1.5pt,
    font=\small
  },
  flowarrow/.style={
    -{Stealth[length=3mm]},
    line width=0.9pt,
    draw=darkaccent
  }
}

\begin{document}

% =================================================
% Титульный лист
% =================================================

\begin{titlepage}
\thispagestyle{empty}
\centering
\vspace*{23mm}

{\Large\textcolor{darkaccent}{\textbf{Чек-лист}}\par}

\vspace{8mm}

{\huge\textbf{Графики показательных функций}\par}

\vspace{3mm}

{\Large\textbf{Метод узловых точек и восстановление формулы}\par}

\vspace{16mm}

{\large ЕГЭ, профильная математика\par}

\vspace{12mm}

{\large\today\par}

\vfill

{\large\textbf{Лёвин Артём Александрович}\par}

\vspace{2mm}

{\large эксклюзивно для Анастасии Павловой\par}

\vspace{20mm}

\begin{tikzpicture}[remember picture,overlay]
  \draw[
    accent!55,
    line width=1.2pt
  ]
    ([xshift=20mm,yshift=19mm]current page.south west)
    .. controls +(35mm,12mm) and +(-35mm,-8mm) ..
    ([xshift=-20mm,yshift=28mm]current page.south east);
\end{tikzpicture}

\end{titlepage}

% =================================================
% Основной чек-лист
% =================================================

\section*{1. Что нужно распознавать сразу}
\sectionline

\begin{itemize}

  \item
  \key{$f(6)$} означает: в найденную формулу функции подставить
  $x=6$ и вычислить значение.

  \item
  \key{«Найдите $x$, если $f(x)=6$»} означает:
  вместо $f(x)$ подставить $6$ и решить полученное уравнение
  относительно $x$.

  \item
  На графике ордината точки --- это значение функции.
  Если точка имеет координаты $(x_0;y_0)$, то
  \[
      f(x_0)=y_0.
  \]

\end{itemize}

\key{Главная развилка:}
сначала выясните, ищете Вы значение функции или аргумент.
Эти формулировки ведут к разным последним шагам.

% -------------------------------------------------

\section*{2. Модель $f(x)=a^x+b$}
\sectionline

Для показательной функции используем ограничения
\[
    a>0,
    \qquad
    a\ne 1.
\]

Параметр $a$ определяет характер показательной части,
а $b$ задаёт вертикальный сдвиг графика.

\begin{itemize}

  \item
  если $a>1$, показательная часть возрастает;

  \item
  если $0<a<1$, показательная часть убывает;

  \item
  при $x=0$ всегда
  \[
      a^0=1,
  \]
  поэтому точка с абсциссой $0$ особенно удобна.

\end{itemize}

% -------------------------------------------------

\section*{3. Узловые точки: рабочая идея}
\sectionline

В задачах по графику удобно искать точки с простыми координатами.
В этом чек-листе будем называть их
\key{узловыми точками}:
обычно это хорошо считываемые точки графика с целыми координатами.

\begin{enumerate}

  \item
  Посчитайте число неизвестных параметров в формуле.

  \item
  Возьмите столько удобных точек,
  сколько неизвестных параметров нужно найти.

  \item
  Предпочитайте точки, где
  \[
      x=0,\qquad x=1,\qquad x=-1
  \]
  или другие небольшие значения.

  \item
  Для каждой точки $(x_0;y_0)$ запишите равенство
  \[
      f(x_0)=y_0.
  \]

  \item
  Подставьте координаты в формулу и решите систему.

\end{enumerate}

\key{Важно:}
если неизвестны два параметра $a$ и $b$,
обычно достаточно двух подходящих точек.
Выписывать все видимые точки графика необязательно.

% -------------------------------------------------

\section*{4. Базовый пример}
\sectionline

Пусть график функции
\[
    f(x)=a^x+b
\]
проходит через точки
\[
    (0;-2)
    \qquad\text{и}\qquad
    (1;-1).
\]

Найдём формулу и вычислим $f(6)$.

Из первой точки:
\[
-2=a^0+b=1+b
\quad\Longrightarrow\quad
b=-3.
\]

Из второй точки:
\[
-1=a^1+b=a-3
\quad\Longrightarrow\quad
a=2.
\]

Следовательно,
\[
    f(x)=2^x-3.
\]

Поэтому
\[
    f(6)=2^6-3=64-3=61.
\]

\begin{center}
\begin{tikzpicture}
\begin{axis}[
    width=10.5cm,
    height=11cm,
    axis lines=middle,
    xmin=-0.5,
    xmax=1.6,
    ymin=-2.5,
    ymax=0.5,
    unit vector ratio=1 1,
    samples=220,
    clip=false,
    xlabel={$x$},
    ylabel={$y$},
    xtick={0,1},
    ytick={-2,-1,0},
    grid=both,
    minor tick num=0,
    grid style={
      line width=.1pt,
      draw=gray!25
    },
    tick label style={
      font=\small
    }
]

\addplot[
  domain=-0.5:1.6,
  very thick,
  accent
]
{2^x-3};

\addplot[
  only marks,
  mark=*,
  mark size=2.3pt,
  darkaccent
]
coordinates {
  (0,-2)
  (1,-1)
};

\node[
  anchor=north east,
  font=\small
]
at (axis cs:-0.08,-2.08)
{$(0;-2)$};

\node[
  anchor=west,
  font=\small,
  xshift=2pt
]
at (axis cs:1,-1)
{$(1;-1)$};

\end{axis}
\end{tikzpicture}
\end{center}

\key{Самопроверка:}
найденная формула должна проходить через обе использованные точки.
Подстановка координат обратно в формулу ---
быстрый способ заметить арифметическую ошибку.

% -------------------------------------------------

\section*{5. Если после подстановки появляется квадрат параметра}
\sectionline

Пусть график
\[
    f(x)=a^x+b
\]
проходит через точки
\[
    (0;-2)
    \qquad\text{и}\qquad
    (2;-1).
\]

Получаем систему
\[
\begin{cases}
-2=a^0+b,\\
-1=a^2+b.
\end{cases}
\]

Из первого уравнения
\[
    b=-3.
\]

Тогда
\[
-1=a^2-3
\quad\Longrightarrow\quad
a^2=2.
\]

Алгебраически уравнение $a^2=2$ имеет два корня,
однако для основания показательной функции действует условие
\[
    a>0.
\]

Поэтому
\[
    a=\sqrt2=2^{1/2}.
\]

Итак,
\[
    f(x)
    =
    \left(2^{1/2}\right)^x-3
    =
    2^{x/2}-3.
\]

Например,
\[
    f(10)
    =
    2^5-3
    =
    29.
\]

\key{Типичная ошибка:}
оставлять значение $a=-\sqrt2$.
Отрицательное число не подходит в качестве основания
показательной функции в рассматриваемой модели.

% -------------------------------------------------

\section*{6. Убывающий график: метод тот же}
\sectionline

Пусть график
\[
    f(x)=a^x+b
\]
проходит через точки
\[
    (-1;-2)
    \qquad\text{и}\qquad
    (0;-3).
\]

Требуется найти $x$, если
\[
    f(x)=12.
\]

Сначала восстанавливаем параметры:
\[
\begin{cases}
-2=a^{-1}+b,\\
-3=a^0+b.
\end{cases}
\]

Из второго уравнения
\[
    b=-4.
\]

Тогда
\[
-2=a^{-1}-4
\quad\Longrightarrow\quad
a^{-1}=2.
\]

Следовательно,
\[
    a=\frac12.
\]

Получаем
\[
    f(x)
    =
    \left(\frac12\right)^x-4.
\]

Теперь используем условие
\[
    f(x)=12:
\]
\[
\left(\frac12\right)^x-4=12.
\]

Отсюда
\[
    \left(\frac12\right)^x=16.
\]

Переходим к основанию $2$:
\[
    2^{-x}=2^4.
\]

Следовательно,
\[
    -x=4,
\]
\[
    x=-4.
\]

\key{Вывод:}
возрастание или убывание графика
не меняет алгоритм восстановления параметров.

% -------------------------------------------------

\section*{7. Когда параметр находится внутри показателя}
\sectionline

Формула может иметь другой вид, например
\[
    f(x)=a^{x+b},
    \qquad
    a>0,
    \qquad
    a\ne1.
\]

Здесь параметр $b$ находится в показателе степени,
поэтому система решается иначе.

Пусть график проходит через точки
\[
    (1;1)
    \qquad\text{и}\qquad
    (3;2).
\]

Тогда
\[
\begin{cases}
1=a^{1+b},\\
2=a^{3+b}.
\end{cases}
\]

При
\[
    a>0,
    \qquad
    a\ne1
\]
равенство
\[
    a^t=1
\]
выполняется только при
\[
    t=0.
\]

Поэтому из первого уравнения:
\[
1+b=0
\quad\Longrightarrow\quad
b=-1.
\]

Во втором уравнении:
\[
2=a^{3-1}=a^2.
\]

Следовательно,
\[
    a=\sqrt2=2^{1/2}.
\]

Получаем
\[
    f(x)
    =
    \left(2^{1/2}\right)^{x-1}
    =
    2^{(x-1)/2}.
\]

Например,
\[
    f(-5)
    =
    2^{(-5-1)/2}
    =
    2^{-3}
    =
    \frac18.
\]

\key{Ключевая проверка:}
всегда смотрите, где именно стоит параметр:
снаружи степени или внутри показателя.
Формулы
\[
    a^x+b
\]
и
\[
    a^{x+b}
\]
требуют разных преобразований.

% -------------------------------------------------

\section*{8. Полезные преобразования степеней}
\sectionline

\begin{align*}
\sqrt a
  &=a^{1/2},
\\[2mm]
\sqrt[n]{a}
  &=a^{1/n},
\\[2mm]
\left(a^m\right)^n
  &=a^{mn},
\\[2mm]
a^{-n}
  &=\frac{1}{a^n},
\\[2mm]
\left(\frac12\right)^x
  &=2^{-x}.
\end{align*}

Если требуется решить уравнение вида
\[
    a^r=c
\]
и число $c$ удаётся представить как степень того же основания,
приводите обе части к одинаковому основанию
и сравнивайте показатели.

% -------------------------------------------------

\section*{9. Типичные ошибки}
\sectionline

\begin{itemize}

  \item
  Перепутать координаты:
  в точке $(x;y)$ первое число подставляется вместо $x$,
  второе --- вместо $f(x)$.

  \item
  Перепутать задачи
  \[
      f(c)
      \qquad\text{и}\qquad
      f(x)=c.
  \]

  \item
  Забыть равенство
  \[
      a^0=1.
  \]

  \item
  Пропустить проверку ограничения
  \[
      a>0
  \]
  после извлечения корня.

  \item
  Ошибиться со знаком отрицательной степени:
  \[
      a^{-1}=\frac1a.
  \]

  \item
  Получить громоздкое или неожиданное число
  и продолжить без проверки исходных точек.
  Необычный результат --- повод быстро перепроверить
  подстановку и арифметику.

  \item
  Считать, что для каждого графика нужен новый способ.
  Метод подстановки координат в заданную модель сохраняется;
  меняется алгебра после подстановки.

\end{itemize}

% =================================================
% Блок-схема алгоритма
% =================================================

\clearpage

\thispagestyle{fancy}

\begin{center}

{\Large\textbf{Алгоритм: восстановление функции по графику}\par}

\vspace{7mm}

\begin{tikzpicture}[node distance=5.2mm and 13mm]

  \node[flowstart] (start)
  {
    Начало: прочитать формулу и понять,
    что требуется найти
  };

  \node[
    flowstep,
    below=of start
  ] (count)
  {
    Посчитать неизвестные параметры
    и выбрать столько же удобных узловых точек
  };

  \node[
    flowstep,
    below=of count
  ] (translate)
  {
    Для каждой точки $(x_0;y_0)$
    записать $f(x_0)=y_0$
  };

  \node[
    flowstep,
    below=of translate
  ] (system)
  {
    Подставить координаты в формулу
    и решить полученную систему
  };

  \node[
    flowstep,
    below=of system
  ] (check)
  {
    Проверить ограничения модели:
    для показательной функции
    $a>0$, $a\ne1$
  };

  \node[
    flowstep,
    below=of check
  ] (restore)
  {
    Записать восстановленную формулу функции
    и проверить использованные точки
  };

  \node[
    flowdecision,
    below=7mm of restore
  ] (goal)
  {
    Что требуется?
  };

  \node[
    flowstep,
    below left=8mm and 6mm of goal,
    text width=5.2cm
  ] (value)
  {
    Найти $f(c)$:
    подставить $x=c$
    и вычислить
  };

  \node[
    flowstep,
    below right=8mm and 6mm of goal,
    text width=5.2cm
  ] (argument)
  {
    Найти $x$ при $f(x)=c$:
    подставить $c$ вместо $f(x)$
    и решить уравнение
  };

  \node[
    flowstart,
    below=28mm of goal
  ] (finish)
  {
    Ответ и финальная проверка
  };

  \draw[flowarrow]
    (start) -- (count);

  \draw[flowarrow]
    (count) -- (translate);

  \draw[flowarrow]
    (translate) -- (system);

  \draw[flowarrow]
    (system) -- (check);

  \draw[flowarrow]
    (check) -- (restore);

  \draw[flowarrow]
    (restore) -- (goal);

  \draw[flowarrow]
    (goal.south west)
    --
    node[
      above left,
      font=\small
    ]
    {\strut $f(c)$}
    (value.north);

  \draw[flowarrow]
    (goal.south east)
    --
    node[
      above right,
      font=\small
    ]
    {\strut $f(x)=c$}
    (argument.north);

  \draw[flowarrow]
    (value.south)
    |-
    ([xshift=-12mm]finish.west)
    --
    (finish.west);

  \draw[flowarrow]
    (argument.south)
    |-
    ([xshift=12mm]finish.east)
    --
    (finish.east);

\end{tikzpicture}

\end{center}

\clearpage

% =================================================
% Тренировочный блок
% =================================================

\section*{Тренировочный блок: 15--16 сентября}
\sectionline

Решайте каждую задачу по одному маршруту:
выбрать удобные точки
$\to$
составить систему
$\to$
найти параметры
$\to$
выполнить требуемое действие.

\subsection*{15.09.2026}

\begin{enumerate}

  \item
  График функции
  \[
      f(x)=a^x+b
  \]
  проходит через точки
  \[
      (0;-1)
      \qquad\text{и}\qquad
      (1;1).
  \]
  Найдите $f(5)$.

  \item
  График функции
  \[
      f(x)=a^x+b
  \]
  проходит через точки
  \[
      (0;-2)
      \qquad\text{и}\qquad
      (2;-1).
  \]
  Найдите $f(8)$.

  \item
  Убывающий график функции
  \[
      f(x)=a^x+b
  \]
  проходит через точки
  \[
      (-1;-2)
      \qquad\text{и}\qquad
      (0;-3).
  \]
  Найдите $x$, если
  \[
      f(x)=12.
  \]

\end{enumerate}

\vspace{5mm}

\subsection*{16.09.2026}

\begin{enumerate}

  \item
  График функции
  \[
      f(x)=a^x+b
  \]
  проходит через точки
  \[
      (0;1)
      \qquad\text{и}\qquad
      (1;3).
  \]
  Найдите $f(4)$.

  \item
  График функции
  \[
      f(x)=a^x+b
  \]
  проходит через точки
  \[
      (0;-4)
      \qquad\text{и}\qquad
      (2;5).
  \]
  Найдите $f(4)$.

  \item
  График функции
  \[
      f(x)=a^{x+b}
  \]
  проходит через точки
  \[
      (1;1)
      \qquad\text{и}\qquad
      (3;4).
  \]
  Найдите $f(-1)$.

\end{enumerate}

% =================================================
% Ответы
% =================================================

\clearpage

\section*{Ответы к тренировочному блоку}
\sectionline

\begin{table}[ht]
\centering
\caption{Ответы для самопроверки}

\begin{tabularx}{\linewidth}{@{} l c X @{}}

\toprule

\textbf{Дата}
&
\textbf{№}
&
\textbf{Ответ}
\\

\midrule

15.09.2026
&
1
&
$f(x)=3^x-2$, поэтому
\[
f(5)=241.
\]
\\

15.09.2026
&
2
&
$f(x)=2^{x/2}-3$, поэтому
\[
f(8)=13.
\]
\\

15.09.2026
&
3
&
$f(x)=\left(\frac12\right)^x-4$, поэтому
\[
x=-4.
\]
\\

\midrule

16.09.2026
&
1
&
$f(x)=3^x$, поэтому
\[
f(4)=81.
\]
\\

16.09.2026
&
2
&
$f(x)=10^{x/2}-5$, поэтому
\[
f(4)=95.
\]
\\

16.09.2026
&
3
&
$f(x)=2^{x-1}$, поэтому
\[
f(-1)=\frac14.
\]
\\

\bottomrule

\end{tabularx}
\end{table}

\vfill

\begin{center}
\textcolor{textgray}{
\small
Перед следующим занятием особенно важно уверенно различать
$f(c)$ и условие $f(x)=c$,
а также без ошибок переводить координаты точки
в равенство $f(x_0)=y_0$.
}
\end{center}

\end{document}